Даны десять вещественных чисел. Найти произведение всех четных чисел.
Даны десять вещественных чисел. Найти среднее арифметическое отрицательных чисел.
Дано целое число $N$ и набор из $N$ вещественных чисел. Вывести сумму и произведение четных чисел из данного набора.
Дано целое число $N$ и набор из $N$ положительных вещественных чисел. Вывести в том же порядке целые части всех чисел из данного набора (как вещественные числа с нулевой дробной частью), а также сумму всех целых частей.
Дано целое число $N$ и набор из $N$ положительных вещественных чисел. Вывести в том же порядке дробные части всех чисел из данного набора (как вещественные числа с нулевой целой частью), а также произведение всех дробных частей.
Дано целое число $N$ и набор из $N$ вещественных чисел. Вывести в том же порядке округленные значения всех чисел из данного набора (как целые числа), а также сумму всех округленных значений.
Дано целое число $N$ и набор из $N$ целых чисел. Вывести в том же порядке все четные числа из данного набора и количество $K$ таких чисел.
Дано целое число $N$ и набор из $N$ целых чисел. Вывести в том же порядке номера всех нечетных чисел из данного набора и количество $K$ таких чисел.
Дано целое число $N$ и набор из $N$ целых чисел. Если в наборе имеются положительные числа, то вывести TRUE; в противном случае вывести FALSE.
Даны целые числа $K$, $N$ и набор из $N$ целых чисел. Если в наборе имеются числа, меньшие $K$, то вывести TRUE; в противном случае вывести FALSE.
Дан набор ненулевых целых чисел; признак его завершения --- число $0$. Вывести количество отрицательных нечетных чисел в наборе.
Дан набор ненулевых целых чисел; признак его завершения --- число $0$. Вывести сумму всех положительных четных чисел из данного набора. Если требуемые числа в наборе отсутствуют, то вывести $0$.
Дано целое число $K$ и набор ненулевых целых чисел; признак его завершения --- число $0$. Вывести количество чисел в наборе, меньших $K$.
Дано целое число $K$ и набор ненулевых целых чисел; признак его завершения --- число $0$. Вывести номер первого числа в наборе, большего $K$. Если таких чисел нет, то вывести $0$.
Дано целое число $K$ и набор ненулевых целых чисел; признак его завершения --- число $0$. Вывести номер последнего числа в наборе, большего $K$. Если таких чисел нет, то вывести $0$.
Дано вещественное число $B$, целое число $N$ и набор из $N$ вещественных чисел, упорядоченных по возрастанию. Вывести элементы набора вместе с числом $B$, сохраняя упорядоченность выводимых чисел.
Дано целое число $N$ и набор из $N$ целых чисел, упорядоченный по возрастанию. Данный набор может содержать одинаковые элементы. Вывести в том же порядке все различные элементы данного набора.
Дано целое число $N$ ($> 1$) и набор из $N$ целых чисел. Вывести те элементы в наборе, которые меньше своего левого соседа, и количество $K$ таких элементов.
Дано целое число $N$ ($> 1$) и набор из $N$ целых чисел. Вывести те элементы в наборе, которые меньше своего правого соседа, и количество $K$ таких элементов.
Дано целое число $N$ ($> 1$) и набор из $N$ вещественных чисел. Проверить, образует ли данный набор возрастающую последовательность. Если образует, то вывести TRUE, если нет --- вывести FALSE.
Дано целое число $N$ ($> 1$) и набор из $N$ вещественных чисел. Проверить, образует ли данный набор убывающую последовательность. Если образует, то вывести TRUE, если нет --- вывести FALSE.
Дано целое число $N$ ($> 1$) и набор из $N$ вещественных чисел. Если данный набор образует убывающую последовательность, то вывести $0$; в противном случае вывести номер первого числа, нарушающего закономерность.
Дано целое число $N$ ($> 2$) и набор из $N$ вещественных чисел. Набор называется пилообразным, если каждый его внутренний элемент либо больше, либо меньше обоих своих соседей (то есть является «зубцом»). Если данный набор является пилообразным, то вывести $0$; в противном случае вывести номер первого элемента, не являющегося зубцом.
Посчитать четные и нечетные цифры введенного натурального числа. Например, если введено число $34560$, то у него $3$ четные цифры ($4$, $6$ и $0$) и $2$ нечетные ($3$ и $5$).
Дано натуральное число $N$. Найти наибольшую и наименьшую цифры числа.
Дано натуральное число $N$. Найти сумму и произведение цифр числа.
Сформировать из введенного натурального числа обратное по порядку входящих в него цифр и вывести на экран. Например, если введено число $3486$, то надо сформировать и вывести число $6843$.
В бригаде, работающей на уборке сена, имеется $N$ косилок. Первая из них работала $M$ часов, а каждая следующая на $10$ минут больше, чем предыдущая. Сколько часов проработала вся бригада?
Составьте программу вывода на экран всех простых чисел, не превосходящих заданного $N$. Простым называется натуральное число больше единицы, имеющее только два делителя: единицу и само это число.
В программе генерируется случайное целое число от $0$ до $100$. Пользователь должен его отгадать не более чем за $10$ попыток. После каждой неудачной попытки должно сообщаться больше или меньше введенное пользователем число, чем то, что загадано. Если за $10$ попыток число не отгадано, то вывести загаданное число.
